\documentclass[11pt]{article}
\usepackage[a4paper,margin=26mm]{geometry}
\usepackage{amsmath,amssymb,mathtools}
\usepackage{unicode-math}
\setmainfont{Latin Modern Roman}
\setmathfont{Latin Modern Math}
\usepackage{xurl,hyperref}
\hypersetup{hidelinks,pdftitle={Problem statement}}
\newcommand{\sref}[2]{\hyperlink{source:#1:#2}{[S#2]}}
\newcommand{\cataloguescope}[1]{#1}
\begin{document}
\section{Numerical invariants — Question 13.5}\label{statement}

\subsection{Definitions and mathematical statement}

\paragraph{Shared notation.}$\mathbb N=\{1,2,\ldots\}$, and $\mathbb Z,\mathbb Q,\mathbb R,\mathbb C$ denote the integers, rationals, reals and complex numbers. Any other conventions are specified locally.


A codimension-one foliation $\mathcal F$ decomposes a $3$-manifold $M$ into surface leaves, locally as the horizontal slices of a product chart. In the terminology of the source, it is taut if a closed curve transverse to $\mathcal F$ meets every leaf. A minimal foliation means that every leaf is dense; a minimal surface means a surface with vanishing mean curvature for the specified Riemannian metric. These two uses of minimal are distinct. Write $\widetilde M$ for the universal cover of $M$ and $\widetilde{\mathcal F}$ for the full lifted foliation. The leaf space $L_{\mathcal F}=\widetilde M/\widetilde{\mathcal F}$ has its quotient topology, and can be non-Hausdorff. A taut foliation is $\mathbb R$-covered if this space is homeomorphic to $\mathbb R$. Its deck action is the holonomy representation $\rho_H:\pi_1(M)\longrightarrow\operatorname{Homeo}(L_{\mathcal F})$. Branching means nonseparated leaves in this space; one-sided branching occurs in only one transverse direction, and two-sided branching in both. For real cubical singular cycles $C=\sum_j a_j\sigma_j$ use $\sum_j|a_j|$. Define a foliated cubical analogue by restricting to cubes whose induced foliation has affine level-set form, and compare with the cubical norm using unrestricted representing cycles.
\paragraph{Statement recorded in the supplied source.}
Develop the foliated Gromov norm using cubical chains and determine whether for every foliation $\mathcal F$ on a hyperbolic $3$-manifold its value on $[M]$ is strictly greater than the usual cubical Gromov norm of $[M]$. Include $\mathbb R$-covered foliations. Both sides must use the same cubical chain convention and normalization. \sref{statement}{2}
\subsection{Short English statement}
Define the cubical foliated norm and test strict inequality over the ordinary cubical fundamental-class norm, including real-covered foliations.
\subsection{Sources}
\begin{description}
\item[\hypertarget{source:statement:1}{S1}] ULAM.AI and the UnsolvedMath Contributors, \emph{UnsolvedMath: A Curated Collection of Open Mathematics Problems}, catalogue entry. CC BY 4.0. \url{https://huggingface.co/datasets/ulamai/UnsolvedMath}.
\item[\hypertarget{source:statement:2}{S2}] Danny Calegari, \emph{Problems in foliations and laminations of 3-manifolds}, version 0.78 (2002), Question 13.5; Definitions 1.1--1.10 and local definitions and remarks. \url{https://arxiv.org/abs/math/0209081v1}.
\end{description}
\label{end:statement}
\end{document}
